\documentclass[11pt]{article}

\usepackage[margin=0.85in]{geometry}
\usepackage{amsmath,amssymb,amsthm,mathtools}
\usepackage{enumitem}
\usepackage{microtype}
\usepackage{xcolor}
\usepackage{hyperref}

\hypersetup{
  colorlinks=true,
  linkcolor=blue,
  urlcolor=blue,
  citecolor=blue
}

\newtheorem{theorem}{Theorem}
\newcommand{\R}{\mathbb R}
\newcommand{\cH}{\mathcal H}
\newcommand{\dd}{\,\mathrm d}
\newcommand{\ee}{\mathrm e}

\title{Davies' Ratio Limit for Locally Perturbed Cylindrical Tubes\\
via the Zero-Threshold Spectral Measure}
\author{Literature-implied answer packet for arXiv:1105.0842}
\date{August 17, 2026}

\begin{document}
\maketitle

\section{Status and source question}

\textbf{Status: literature-implied positive answer (full for the two classes
asked about).}  This packet records an agent-identified consequence of an
existing cylindrical-threshold spectral theorem.  It is not presented as a
new theorem of the supporting authors.

Grillo, Kova\v{r}\'ik, and Pinchover ask in Appendix~C, Remark~C.2, on PDF
page~22 of the current arXiv PDF~\cite{GKP} whether Davies' conjecture holds

\begin{quote}
for operators \(L\) satisfying the assumptions of Theorem~3.19 and for the
Laplacian on a twisted tube.
\end{quote}

Here Theorem~3.19 concerns the nonnegative self-adjoint Friedrichs realization
of a real uniformly elliptic operator
\[
 L=-\nabla\!\cdot(a\nabla)+V
\]
on the straight tube \(\Omega_0=\omega\times\R\), assumed subcritical and
equal to \(-\Delta-E_1\) outside a compact longitudinal region.  The twisted
case is the shifted Dirichlet Laplacian
\(H=-\Delta^D_\Omega-E_1\), where the twisting is compactly supported.
The source proves, for either class and each fixed interior point \(x\),
\begin{equation}\label{eq:diag}
 c_x t^{-3/2}\leq \ee^{-tH}(x,x)\leq C_x t^{-3/2}
 \qquad(t\geq t_x).
\end{equation}

The decisive later input is Christiansen--Datchev's black-box spectral theory
for compactly supported perturbations of cylindrical ends~\cite{CD}.  Their
Sections~2.2--2.5 cover multiple ends and Dirichlet waveguides; Lemmas~2.4,
2.5, and~2.6 give threshold regularity, the generalized-eigenfunction spectral
measure, and the form of a threshold resonance.  The supporting paper does
not mention Davies' conjecture or identify this application.

\section{Identification and result}

Both source classes fit the Christiansen--Datchev framework with
\[
 H_Y=-\Delta^D_\omega-E_1\geq0
\]
on each cylindrical end.  The zero eigenvalue of \(H_Y\) has finite
multiplicity \(q\) (in the two-ended tube, \(q=2\)), and the next transverse
threshold \(\nu_1\) is positive.  Let
\(\Phi_1(\lambda),\ldots,\Phi_q(\lambda)\) be the generalized eigenfunctions
associated with the zero channels, normalized as in~\cite[Section~2.3]{CD}.

\begin{theorem}\label{thm:main}
For every operator \(H=L\) in source Theorem~3.19, and for the shifted
Dirichlet Laplacian \(H=-\Delta^D_\Omega-E_1\) of the locally twisted tube,
there are locally continuous functions (with the elliptic regularity supplied
by the coefficients)
\[
 \Psi_j=\partial_\lambda\Phi_j(0),\qquad 1\leq j\leq q,
\]
such that, for every fixed pair of interior points \(x,y\),
\begin{equation}\label{eq:asymptotic}
 \ee^{-tH}(x,y)
 =t^{-3/2}B(x,y)+O_{x,y}(t^{-2}),
 \qquad
 B(x,y)=\frac{1}{8\sqrt\pi}
 \sum_{j=1}^q\Psi_j(x)\overline{\Psi_j(y)}.
\end{equation}
Moreover \(B(x,y)>0\) for all \(x,y\).  Consequently, for arbitrary fixed
reference points \(x_0,y_0\),
\begin{equation}\label{eq:ratio}
 \lim_{t\to\infty}
 \frac{\ee^{-tH}(x,y)}{\ee^{-tH}(x_0,y_0)}
 =\frac{B(x,y)}{B(x_0,y_0)}>0.
\end{equation}
The same ratio limit holds for the unshifted twisted-tube Dirichlet heat
kernel.  Thus Davies' conjecture has a positive answer in both cases singled
out in Remark~C.2.
\end{theorem}

\begin{proof}
We give the derivation because the implication is not stated in the supporting
paper.

\smallskip
\noindent\emph{1. Exclusion of zero point spectrum.}
A zero eigenfunction would contribute a nondecaying squared-modulus term to
the diagonal heat kernel at every point where it is nonzero.  This contradicts
the upper bound in~\eqref{eq:diag}.  Hence zero is not an eigenvalue of \(H\).

\smallskip
\noindent\emph{2. Exclusion of a zero-threshold resonance.}
Let \(R(\lambda)=(H-\lambda^2)^{-1}\), with \(R(\lambda)\) and
\(R(-\lambda)\) denoting the two physical boundary values for \(\lambda>0\).
Below the first positive transverse threshold, Christiansen--Datchev
Lemma~2.5 says, after compact localization,
\begin{equation}\label{eq:specmeasure}
 \frac1i\bigl(R(\lambda)-R(-\lambda)\bigr)
 =\frac1{2\lambda}\sum_{j=1}^q
   \Phi_j(\lambda)\otimes\Phi_j(\lambda),
 \qquad 0<\lambda<\nu_1,
\end{equation}
apart from the point-spectrum projection, which has already been excluded at
zero.  If some \(\Phi_j(0)\) were nonzero, choose \(x\) with
\(\Phi_j(0,x)\neq0\).  Threshold continuity from their Lemma~2.4 and the
positive diagonal spectral density in~\eqref{eq:specmeasure} would imply
\[
 \ee^{-tH}(x,x)
 \geq c\int_0^\varepsilon \ee^{-t\lambda^2}\dd\lambda
 \geq c' t^{-1/2}
\]
for all sufficiently large \(t\), contradicting~\eqref{eq:diag}.  Therefore
\begin{equation}\label{eq:nores}
 \Phi_j(0)=0\qquad(1\leq j\leq q).
\end{equation}
By Christiansen--Datchev Lemma~2.6 and the discussion following it, this is
exactly the absence of a zero-threshold resonance, so the cutoff resolvent is
regular at zero.

\smallskip
\noindent\emph{3. The exact heat-kernel asymptotic.}
The generalized eigenfunctions are meromorphic in the threshold coordinate
and regular at the physical threshold; at zero that coordinate is simply
\(\lambda\).  Local elliptic regularity lets this expansion be evaluated
pointwise on compact subsets.  Hence~\eqref{eq:nores} gives
\begin{equation}\label{eq:phiexp}
 \Phi_j(\lambda)=\lambda\Psi_j+O(\lambda^2).
\end{equation}
Choose \(0<\delta<\nu_1\) so small that the meromorphically continued cutoff
resolvent has no pole for \(|\lambda|<\delta\).  Stone's formula and
\eqref{eq:specmeasure} give the low-energy kernel
\begin{equation}\label{eq:stone}
 K_{<\delta^2}(t;x,y)
 =\frac1{2\pi}\int_0^\delta \ee^{-t\lambda^2}
   \sum_{j=1}^q
   \Phi_j(\lambda,x)\overline{\Phi_j(\lambda,y)}\dd\lambda.
\end{equation}
Indeed, the factor \(2\lambda\) from \(\dd(\lambda^2)\) cancels the
\(2\lambda\) in~\eqref{eq:specmeasure}.  Substituting~\eqref{eq:phiexp} and
using
\[
 \int_0^\infty\ee^{-t\lambda^2}\lambda^2\dd\lambda
 =\frac{\sqrt\pi}{4t^{3/2}}
\]
yields~\eqref{eq:asymptotic} for the low-energy part, with an
\(O_{x,y}(t^{-2})\) remainder.

For completeness, the remaining spectrum is exponentially negligible
pointwise.  If \(E_H\) is the spectral resolution, insert the time-\(1/2\)
heat kernel on both sides of
\(E_H([\delta^2,\infty))\ee^{-(t-1)H}\).  Since
\(\|\ee^{-H/2}\delta_x\|_2^2=\ee^{-H}(x,x)<\infty\), Cauchy--Schwarz gives
\[
 \bigl|K_{\geq\delta^2}(t;x,y)\bigr|
 \leq \ee^{-(t-1)\delta^2}
 \sqrt{\ee^{-H}(x,x)\ee^{-H}(y,y)}.
\]
This also absorbs any positive embedded eigenvalues away from zero.

\smallskip
\noindent\emph{4. Strict positivity.}
Taking \(x=y\) in~\eqref{eq:asymptotic} and using the lower half of
\eqref{eq:diag} gives \(B(x,x)>0\) at every interior point.  Fixed-point
parabolic Harnack comparison gives, for each fixed \(x,y\), a positive lower
comparison between \(\ee^{-tH}(x,y)\) and a diagonal kernel at a boundedly
shifted time.  In the twisted-tube discussion the source states this
pointwise equivalence explicitly, citing Davies' Theorem~10
\cite[PDF p.~2]{GKP}.  Passing to~\eqref{eq:asymptotic} gives
\(B(x,y)>0\).  Dividing the two asymptotics proves~\eqref{eq:ratio}.

Finally, if \(P=-\Delta^D_\Omega\), then
\(\ee^{-tP}=\ee^{-E_1t}\ee^{-t(P-E_1)}\); the scalar exponential cancels
from the ratio.
\end{proof}

\section{Scope, provenance, and review points}

This settles exactly the self-adjoint uniformly elliptic class in source
Theorem~3.19 and its locally twisted Dirichlet tube.  It does not address
arbitrary nonsymmetric parabolic operators or Davies' conjecture on general
noncompact manifolds.

The run's registry, solution, attempt, and proof-gap indexes were searched for
the source id and the core Davies/twisted-tube/ratio-limit terms.  Bounded
direct searches covered the exact source phrase, the source title, ``Davies
conjecture twisted tube,'' and later heat/spectral asymptotics for cylindrical
waveguides.  No paper explicitly announcing this answer was found.  The
supporting paper supplies the decisive existing spectral identity and
threshold regularity, but does not cite the source or discuss this ratio
limit.  The classification is therefore \emph{literature-implied answer}, not
\emph{literature already answered}, and no originality claim is made.

The main human-review points are: (i) the black-box fit for a compactly
supported coefficient perturbation after shifting by \(E_1\); (ii) the
localized Stone normalization in~\eqref{eq:stone}; (iii) the pointwise
passage, justified by local elliptic regularity and the time-one smoothing
estimate; and (iv) the Harnack step that upgrades a nonzero diagonal
coefficient to strict positivity for every ordered pair.

\begin{thebibliography}{9}

\bibitem{GKP}
G. Grillo, H. Kova\v{r}\'ik, and Y. Pinchover,
\emph{Sharp two-sided heat kernel estimates of twisted tubes and
applications}, arXiv:\href{https://arxiv.org/abs/1105.0842}{1105.0842}.

\bibitem{CD}
T. J. Christiansen and K. Datchev,
\emph{Wave asymptotics for waveguides and manifolds with infinite
cylindrical ends},
arXiv:\href{https://arxiv.org/abs/1705.08972}{1705.08972}.

\end{thebibliography}

\end{document}
